\documentclass[a4paper]{article}
% Español (México) con XeLaTeX: fontspec para fuentes Unicode, polyglossia
% para la división silábica y los nombres del documento en la variante mexicana.
\usepackage{fontspec}
\usepackage{polyglossia}
\setdefaultlanguage[variant=mexican]{spanish}
% Los nombres chinos van como «pinyin (original)»: xeCJK da a los caracteres
% Han su propia fuente; sin ella XeLaTeX los omite con un aviso.
% FandolSong no cubre algunos caracteres de nombres (U+73FA); WenQuanYi Zen Hei sí.
\usepackage[AutoFallBack=true]{xeCJK}
\setCJKmainfont{FandolSong-Regular.otf}
\setCJKfallbackfamilyfont{\CJKrmdefault}{WenQuanYi Zen Hei}
% polyglossia no traduce los nombres de listings.
\AtBeginDocument{\providecommand{\lstlistingname}{}\renewcommand{\lstlistingname}{Listado}%
  \providecommand{\lstlistlistingname}{}\renewcommand{\lstlistlistingname}{Índice de listados}}
\usepackage{amsmath, amssymb}
\usepackage{graphicx}
\usepackage[margin=2.5cm]{geometry}
\usepackage[most]{tcolorbox}
\usepackage{etoolbox}
\usepackage{listings}
\usepackage{hyperref}
\usepackage{booktabs}
\usepackage{float}
\newtcolorbox{knowledgebox}[1]{enhanced, colback=blue!5!white, colframe=blue!75!black, colbacktitle=blue!75!black, coltitle=white, fonttitle=\bfseries, title={#1}, attach boxed title to top left={yshift=-2mm, xshift=2mm}, boxrule=1pt, sharp corners}
\newtcolorbox{importantbox}[1]{enhanced, colback=yellow!10!white, colframe=yellow!80!black, colbacktitle=yellow!80!black, coltitle=black, fonttitle=\bfseries, title={#1}, sharp corners}
\newtcolorbox{warningbox}[1]{enhanced, colback=red!5!white, colframe=red!75!black, colbacktitle=red!75!black, coltitle=white, fonttitle=\bfseries, title={#1}, sharp corners}
\lstset{language=Python, basicstyle=\ttfamily\small, keywordstyle=\color{blue}, stringstyle=\color{red!60!black}, commentstyle=\color{green!60!black}, breaklines=true, frame=single, numbers=left, numberstyle=\tiny\color{gray}, captionpos=b, extendedchars=true}
\newcommand{\notetitle}{Complemento a los detalles del entrenamiento de SFT: aprendizaje supervisado sin necesidad de Decoding}
\newcommand{\noteauthors}{Elaboradas a partir de materiales públicos del curso}
\newcommand{\notedate}{2025}
\newcommand{\videochannel}{Wudaokou Nash (五道口纳什)}
\newcommand{\videopublishdate}{2025}
\newcommand{\videoduration}{aproximadamente 25 minutos}
\newcommand{\videourl}{}
\newcommand{\videocoverpath}{cover.jpg}
\newcommand{\repourl}{https://github.com/hqhq1025/ai-course-notes}

% Footer with repo info
\usepackage{fancyhdr}
\pagestyle{fancy}
\fancyhf{}
\fancyfoot[C]{\thepage}
\fancyfoot[R]{\footnotesize\color{gray}\href{\repourl}{AI Course Notes}}
\renewcommand{\headrulewidth}{0pt}
\renewcommand{\footrulewidth}{0pt}

\begin{document}
\begin{titlepage}\centering
{\Large Notas del curso\par}\vspace{1.2cm}{\huge\bfseries \notetitle\par}\vspace{0.8cm}{\large \noteauthors\par}\vspace{0.3cm}{\large \notedate\par}\vspace{1.2cm}
\ifdefempty{\videocoverpath}{}{\includegraphics[width=0.82\textwidth,height=0.45\textheight,keepaspectratio]{\videocoverpath}\par}
\vfill
\begin{tcolorbox}[width=0.9\textwidth, colback=black!2!white, colframe=black!60, sharp corners]
\textbf{Autor o canal del video}: \videochannel\par\textbf{Fecha de publicación}: \videopublishdate\par\textbf{Duración del video}: \videoduration\par
\ifdefempty{\videourl}{}{\textbf{Enlace del video}: \href{\videourl}{\nolinkurl{\videourl}}\par}\par
\textbf{Página del proyecto}: \href{\repourl}{\nolinkurl{\repourl}}\par
\end{tcolorbox}
\end{titlepage}
\tableofcontents\newpage

\section{Introducción}
Esta entrega complementa más detalles del SFT: el entrenamiento de SFT no necesita decoding, así como el cálculo de la Cross-Entropy Loss.

\begin{importantbox}{El entrenamiento de SFT no necesita Decoding}
El SFT es aprendizaje supervisado, y ya conocemos el target response. Durante el entrenamiento el modelo \textbf{no necesita} generar token por token (decoding), sino que calcula de una sola vez y en paralelo la loss de todas las posiciones. Esto es esencialmente distinto de la auto-regressive generation durante la inferencia.
\end{importantbox}

\section{Cross-Entropy Loss en detalle}
Para modelos decode-only como GPT:
\begin{itemize}
  \item Se recibe la secuencia completa (prompt + response) como entrada
  \item El modelo produce en paralelo los logits de todas las posiciones
  \item Solo se calcula la cross-entropy loss en las posiciones de los tokens de response
  \item La parte del prompt se enmascara mediante el loss mask
\end{itemize}

\section{Training vs. Inference}
\begin{warningbox}{Diferencia clave}
\begin{itemize}
  \item \textbf{Training}: Teacher Forcing, calcula en paralelo todas las posiciones
  \item \textbf{Inference}: Auto-regressive, genera token por token
\end{itemize}
\end{warningbox}

\subsection{Resumen de la sección}
La eficiencia del SFT proviene de Teacher Forcing: calcula en paralelo el loss de todas las posiciones, sin necesidad de decodificar token por token.

\newpage
\section{Origen de la eficiencia: por qué entrenamiento e inferencia difieren tanto}
Muchos principiantes, la primera vez que se topan con el SFT, se preguntan: si es el mismo modelo de lenguaje, ¿por qué el entrenamiento puede calcular en paralelo todas las posiciones, mientras que la inferencia debe generar un token a la vez? Esta es precisamente la diferencia esencial entre Teacher Forcing y Auto-regressive decoding.

\begin{importantbox}{La diferencia en una frase}
En la etapa de entrenamiento se conoce la secuencia objetivo completa, por lo que se puede calcular en paralelo la loss de todas las posiciones de la response; en la etapa de inferencia no se conoce el siguiente token, así que solo se puede generar paso a paso. Esta es la razón de fondo por la que el SFT suele ser mucho más barato que el rollout de RL.
\end{importantbox}

\subsection{Resumen de la sección}
Entender el Teacher Forcing no es solo entender una técnica de entrenamiento, sino entender por qué el costo de sistema de la etapa de aprendizaje supervisado y el de la etapa de rollout en línea son tan distintos.

\newpage
\section{Síntesis y ampliación}
\begin{enumerate}
  \item El entrenamiento de SFT no necesita decoding (Teacher Forcing)
  \item El Cross-Entropy Loss se calcula solo sobre la parte de response
  \item Esta es la razón fundamental por la que el SFT es mucho más rápido que el entrenamiento con RL
\end{enumerate}
\end{document}
